\section{Background and related work}\label{sec:background}
\paragraph{Clipped policy optimization.}
Let $q$ be a prompt, $o=(o_1,\ldots,o_T)$ a sampled response, and $h_t=(q,o_{<t})$ its token context. Rollouts are generated by a fixed policy $\pi_{\mathrm{old}}$, and $\pi_\theta$ is the policy being updated. Define the token probability ratio and clipped surrogate as
\begin{align}
 s_t(\theta)&=\frac{\pi_\theta(o_t\mid h_t)}{\pi_{\mathrm{old}}(o_t\mid h_t)},\\
 f_t(\theta,A_t)&=\min\!\left\{s_t(\theta)A_t,\,
 \operatorname{clip}(s_t(\theta),1-\varepsilon,1+\varepsilon)A_t\right\}.
 \label{eq:clipped}
\end{align}
The response-averaged PPO objective is
\begin{equation}
 J_{\mathrm{clip}}(\theta)=\E_{q,\,o\sim\pi_{\mathrm{old}}(\cdot\mid q)}
 \left[\frac{1}{T}\sum_{t=1}^{T}f_t(\theta,A_t)\right].
 \label{eq:ppo}
\end{equation}
PPO commonly obtains $A_t$ from a learned value function using generalized advantage estimation~\citep{schulman2018highdimensionalcontinuouscontrolusing}. Critic-free methods instead construct $A_t$ from observed rewards. All advantages are held fixed while optimizing \cref{eq:ppo}.

\paragraph{Baselines and normalization are separate choices.}
For $k$ responses to the same prompt, write $r_i=r(q,o^{(i)})$, $\bar r_q=k^{-1}\sum_i r_i$, and $s_q^2=k^{-1}\sum_i(r_i-\bar r_q)^2$. ReMax subtracts the reward of a greedy response. RLOO uses
\begin{equation}
 A_i^{\mathrm{RLOO}}=r_i-\frac{1}{k-1}\sum_{j\ne i}r_j,
 \qquad k>1.
 \label{eq:rloo}
\end{equation}
GRPO uses the group-standardized reward
\begin{equation}
 A_i^{\mathrm{GRPO}}=\frac{r_i-\bar r_q}{s_q+\eta},
 \qquad \eta>0.
 \label{eq:grpo}
\end{equation}
Thus, the group mean controls centering, while $s_q$ controls the relative scale of different prompts. Each group receives its own scale, estimated from $k$ responses. Our method retains the option of a group baseline while sharing the normalization scale across the rollout batch. \Cref{appendix:normalization} analyzes the finite-sample behavior of group standardization and its large-sample limit.

\paragraph{Normalization in RL.}
Advantage normalization is an established implementation choice in policy optimization~\citep{andrychowicz2020matters}. Our focus is its use in critic-free language-model training and its interaction with prompt baselines, KL regularization, and response sampling. Related studies examine RL implementation choices~\citep{liu2025part}, compute scaling~\citep{khatri2025art}, and batch-wise normalization in length-constrained reasoning~\citep{liu2025dler}. Together, these studies emphasize the practical importance of normalization and loss design across language-model RL settings.
